\documentclass[12pt]{article}

\usepackage{sbc-template}
\usepackage[brazil]{babel}
\usepackage{graphicx}
\usepackage{times}
\usepackage{url}

\sloppy

\title{Ideias para adiar o fim do mundo: Uma Leitura Pessoal sobre a Resistência e a Reconexão com a Terra em Ailton Krenak}

\author{Yago Melo da Costa}

\begin{document}
\maketitle

\begin{resumo}
Este artigo de opinião apresenta uma análise subjetiva da obra "Ideias para adiar o fim do mundo", de Ailton Krenak, a partir da leitura do autor no ano de 2026. O texto explora a crítica do autor à humanidade dissociada da natureza, a denúncia do extrativismo capitalista e a proposta de resistência através das subjetividades, da memória ancestral e do sonho. Através de uma reflexão pessoal, a obra é interpretada como um alerta existencial sobre o Antropoceno e um convite à reinvenção, dialogando com os artigos anteriores do leitor sobre as obras de Saramago e Kafka, num percurso que vai da alienação do indivíduo à necessidade de pertencimento coletivo à Terra.
\end{resumo}

\begin{abstract}
This opinion article presents a subjective analysis of Ailton Krenak's "Ideias para adiar o fim do mundo", based on the author's reading in 2026. The text explores the author's critique of humanity dissociated from nature, the denunciation of extractivist capitalism, and the proposal for resistance through subjectivities, ancestral memory, and dreams. Through personal reflection, the work is interpreted as an existential warning about the Anthropocene and an invitation to reinvention, dialoguing with the reader's previous articles on the works of Saramago and Kafka, in a journey from individual alienation to the need for collective belonging to the Earth.
\end{abstract}

\section{Introdução}

\noindent
\textbf{Edição utilizada:} Todas as citações e referências a números de página neste artigo referem-se à 1ª edição publicada pela Companhia das Letras \cite{krenak2019}, que reúne três textos: a palestra que dá título ao livro, "Do sonho e da terra" e "A humanidade que pensamos ser". Foi a única edição consultada para esta leitura.

Após ter me aventurado pelas angústias de Gregor Samsa em "A Metamorfose" e pela busca existencial do homem anônimo em "O Conto da Ilha Desconhecida", cheguei a Ailton Krenak. Esperava um ensaio político, mas me deparei com um grito visceral, uma prece filosófica e um chamado urgente. "Ideias para adiar o fim do mundo" não é um livro de receitas; é uma constatação dolorosa e uma poética da resistência. A princípio, sua abordagem é radical, porém percebi o quão necessária à crítica a um progresso ou desenvolvimento sustentável, visto que a causa do desastre ambiente não é tratada. O modo de produção que vivemos pouco se preocupa com a manutenção do mundo mas apenas com o aumento dos lucros aos poucos através da extração dos bens gerados pela natureza. À medida que fui acompanhando sua fala, fui tomado por uma inquietação profunda, que logo se transformou em um misto de indignação e esperança. Este artigo é uma tentativa de organizar o choque dessa leitura, conectando a visão de Krenak aos meus percursos literários anteriores e ao meu próprio lugar no mundo.

\section{Sobre a Obra e a Primeira Impressão}

O livro começa com uma cena memorável: Krenak no quintal de casa, sendo convidado para uma palestra sobre desenvolvimento sustentável na UnB, e respondendo, no impulso, com o título que daria origem à obra. Essa anedota inicial já carrega toda a potência do livro. Krenak desmonta o mito da sustentabilidade com uma ironia cortante, revelando que ela é, muitas vezes, um "assalto à nossa ideia de natureza" \cite[p. 10]{krenak2019}.

Minha primeira impressão foi de estranhamento. Krenak questiona a própria base do que chamamos de humanidade. Ele nos pergunta: "Será que ela não está na base de muitas das escolhas erradas que fizemos, justificando o uso da violência?" \cite[p. 7]{krenak2019}. Me vi confrontado com uma ideia que parece óbvia, mas que nos é sistematicamente ensinada a ignorar: a de que somos parte da natureza e não seus senhores. A "humanidade que pensamos ser" é uma abstração que nos descolou da terra, e essa foi a primeira grande chave para começar a compreender a profundidade do que ele estava dizendo.

\section{A Crítica ao "Clube da Humanidade" e ao Consumismo}

O que mais me impactou foi a forma como Krenak expõe a falência do nosso modelo civilizatório. Ele nos lembra que a ideia de que "os brancos europeus podiam sair colonizando o resto do mundo" estava baseada na premissa de que "havia uma humanidade esclarecida que precisava ir ao encontro da humanidade obscurecida" \cite[p. 7]{krenak2019}. Essa narrativa de salvação justificou séculos de violência e exploração.

Ele denuncia, ainda, que a modernização jogou milhões de pessoas de seus territórios para centros urbanos, arrancando-as de suas memórias ancestrais e as transformando em meras peças de um "liquidificador chamado humanidade". O autor critica com precisão a transformação do cidadão em consumidor, citando José Mujica: "Não tem gente mais adulada do que um consumidor. São adulados até o ponto de ficarem imbecis, babando" \cite[p. 14]{krenak2019}.

Essa denúncia me lembrou diretamente de Gregor Samsa. Se Kafka mostrou a metamorfose do indivíduo em um inseto inútil ao capital, Krenak mostra a metamorfose da própria humanidade em uma legião de zumbis consumistas, alienados e desconectados da sua própria essência. A diferença é que Krenak não nos deixa apenas na agonia do quarto escuro de Gregor; ele aponta a porta para sairmos dele.

\section{A Natureza como Sujeito e a Dor do Watu}

A parte do livro que mais me tocou foi o relato de Krenak sobre o rio Doce, chamado por seu povo de \emph{Watu}, nosso avô. Krenak conta que, após o rompimento da barragem de rejeitos de mineração, o rio, que sustentou seu povo por gerações, entrou em coma. "Aquele rio que está em coma é também o nosso avô", afirma ele \cite[p. 26]{krenak2019}.

Ele vai além: fala sobre a montanha Takukrak, que possui personalidade, "cara", e que é capaz de indicar se o dia vai ser bom ou se é melhor ficar quieto. A ideia de despersonalizar a natureza, de transformá-la em "recurso", é vista por ele como uma violência ontológica. "Quando despersonalizamos o rio, a montanha, quando tiramos deles os seus sentidos, considerando que isso é atributo exclusivo dos humanos, nós liberamos esses lugares para que se tornem resíduos da atividade industrial e extrativista" \cite[p. 26]{krenak2019}.

Isso ecoou em mim uma tristeza semelhante à que senti ao ler "A Metamorfose". Mas, enquanto no livro de Kafka a tristeza vinha da solidão de um filho rejeitado, aqui a tristeza é amplificada: é a solidão de um planeta sendo rejeitado por seus próprios filhos. A família de Krenak é o rio, a montanha, a terra, e a "família humana" (as corporações, os governos) o rejeitou e o está matando. A angústia é a mesma, mas a escala é cósmica.

\section{Os Paraquedas Coloridos e o Sonho como Disciplina}

Apesar do tom de denúncia, a obra de Krenak não é um beco sem saída. No último capítulo, ele nos presenteia com uma imagem poderosa: a do paraquedas colorido. Krenak questiona por que temos tanto medo de cair, se ao longo da história não fizemos outra coisa senão despencar. Ele nos convida a não eliminar a queda, mas a "inventar e fabricar milhares de paraquedas coloridos, divertidos, inclusive prazerosos" \cite[p. 33]{krenak2019}.

Ele nos fala sobre o sonho como uma "disciplina", um exercício de busca de orientação para a vida, e não um mero devaneio. Ele nos chama a "suspender o céu", a expandir o horizonte existencial, e não a nos contentarmos com a visão rasteira do consumo e do produtivismo. "A minha provocação sobre adiar o fim do mundo é exatamente sempre poder contar mais uma história" \cite[p. 15]{krenak2019}.

Essa ideia do paraquedas colorido ressoou com o tema da "ilha desconhecida" de Saramago. Se Saramago nos convidava a embarcar em direção a um destino incerto, Krenak nos convida a mergulhar de cabeça na queda inevitável do nosso tempo, mas com a leveza e a criatividade de quem se recusa a ser apenas um número, de quem se recusa a ser apenas consumidor e insiste em ser um contador de histórias.

\section{Minha Interpretação Subjetiva e Diálogo com Leituras Anteriores}

Minha percepção central da obra é que Krenak está nos oferecendo um antídoto para duas doenças que já explorei em minhas leituras anteriores. De um lado, há a doença da alienação e da falta de propósito (representada por Gregor Samsa, e pela angústia de viver em um sistema que nos reduz a funções). De outro, está a doença do tédio e da burocratização da existência (representada pelo palácio do rei em Saramago e a negação da existência de uma ilha desconhecida).

Krenak trata o sistema capitalista e a modernidade como a verdadeira "metamorfose": a metamorfose da humanidade em uma abstração que consome a si mesma. O "besouro" de Kafka é o estado terminal do consumidor alienado. Krenak, no entanto, nos dá a ferramenta para reverter essa transformação: a reconexão. Ele nos lembra que o rio é o avô, a montanha é a avó, e nós somos a terra.

A obra me fez ressignificar o conceito de "ilha desconhecida" de Saramago. Para mim, a ilha agora não é apenas um destino exterior para o autoconhecimento, mas é o retorno à nossa própria natureza orgânica. É perceber que somos feitos de água, terra e ar, e que o "desenvolvimento sustentável" é apenas uma cortina de fumaça para esconder o abismo que estamos cavando. A mulher da limpeza de Saramago, que agia com coragem e desprendimento, se torna aqui a figura do sábio indígena que sabe que a terra não é um recurso, mas um parente.

\section{Reflexões Pessoais e Conclusão}

Ler "Ideias para adiar o fim do mundo" foi um exercício de desconstrução pessoal. Como um habitante da cidade, criado sob o signo do consumo e do progresso, me senti confrontado. Foi um livro incômodo, pois me fez enxergar o quanto minha própria vida, minhas escolhas e meu conforto estão diretamente ligados à destruição que ele denuncia.

Ainda assim, a sensação de urgência que ele provoca não me levou ao desespero, mas à ação. Assim como nos livros anteriores, onde a literatura me fez refletir sobre meu lugar no mundo, Krenak me fez querer, na medida do possível, me descolar dessa "abstração civilizatória". A ideia de "viver no prazer aqui na Terra" sem culpa, sem a necessidade de justificar cada ato pela produtividade, é um convite libertador.

Senti a necessidade de buscar meus próprios "paraquedas coloridos": pequenos gestos de reconexão, de valorização da diversidade e de recusa ao consumo desenfreado. A obra de Krenak não me oferece respostas prontas, mas me dá a coragem de continuar caindo, de continuar contando minhas próprias histórias e, quem sabe, de "adiar o fim do mundo" por mais alguns dias.

Ao terminar o livro, compreendi que a jornada de leitura que comecei com Kafka e Saramago chegou a um ponto de encontro. Da metamorfose individual de Gregor à busca pela ilha de Saramago, e agora à advertência coletiva de Krenak, percebo que a literatura é um fio que costura não apenas histórias, mas possibilidades de existência. E a maior delas, ensinada por Krenak, é a de que somos capazes de nos reinventar, não como deuses, mas como parte integrante e pulsante da vida na Terra.

\begin{thebibliography}{99}

\bibitem{krenak2019}
KRENAK, Ailton.
\emph{Ideias para adiar o fim do mundo}.
1ª ed. São Paulo: Companhia das Letras, 2019.
ISBN 978-85-359-3241-6.

\end{thebibliography}

\end{document}